% Shift-JIS CRLF 版
\documentclass[a4paper,12pt]{jsreport}
\usepackage{bm}
\usepackage[dvipdfmx]{graphicx}
\usepackage{ascmac}
\usepackage{amsmath}
\usepackage{float}
\newcommand{\argmax}{\mathop{\rm arg~max}\limits}

\title{類似貢献度に基づく文書検索結果の可視化分析法と結果の評価法}
\author{神奈川大学大学院　理学研究科理学専攻 \\
202370036　尾崎 了祐}
\date{2025年}

\begin{document}
\maketitle
\tableofcontents

\chapter{はじめに}

ビッグデータ技術の進展により, 膨大な文書データから, 有用な情報を利活用するニーズは益々高まっている [1].
本研究では, 与えられた文書クエリの類似検索結果文書集合を可視化する課題に焦点を当てる. 
ところが, 著者らの知る限り, この課題に関する先行研究は少数で~\cite{oda04}, あまり研究が進んでいない状況である. 
本論文では, 検索結果の各文書をノードと見なし, それぞれを類似貢献度ベクトルに変換する. 
そして, k-最近傍グラフを構築した可視化において, 検索結果文書がクエリに対し, 
どのような説明語で類似するかを付与するとともに, グループ分けして配色する新手法を提案する. 
livedoor ニュースコーパスを用いた評価実験では, 可視化例に基づく定性評価とともに, 
説明語となる単語数の定量評価により, 提案法の有効性を検証する.
さらに, 本論文では, 各種可視化法の有効性を, エッジ長の変動係数を用いて定量評価する方法を提案する. 
livedoorニュースコーパスを用いた評価実験では, ４つの手法でグラフを可視化し, 
定量評価と可視化例に基づく定性評価を比較することで, 定量評価法の妥当性を検証する.

本論文の構成は以下の通りである. 
2章では, 先行研究を述べる.
3章では, 提案分析法について詳細に述べる. 
4章では, 実験の詳細設定と, 実験で得られた結果を示す. 
5章では, 先行研究での結果と提案法での結果を比較し, 提案法の有効性を示す.
6章では，類似貢献度に基づく文書検索結果の可視化評価法ついて詳細に述べ，実験で得られた結果を示す．
7章では, これまでの成果を応用して開発した可視化ブラウジングツールを紹介する.
最後に, 本研究のまとめと, 今後の展望を述べる. 

\newpage

\chapter{先行研究}
大規模文書検索結果のクラスタリングと可視化を組み合わせている先行研究として,~\cite{oda04}が挙げられる.
当研究は, 検索結果文書全体を対象として, 文書間の類似度を考慮した可視化結果を生成している点で共通しているが,
文書のクラスタリング手法として, 文書間の距離ベクトルを最長距離法で算出する階層型クラスタリングを採用しており,
検索結果文書を3次元空間上で可視化している.
一方で, 本研究では文書間の類似ベクトルをコサイン類似度で算出し,
そのベクトルをクエリを考慮した類似貢献度ベクトルに変換する,非階層型クラスタリングを採用している. 
検索結果文書も, 2次元上で可視化している点で異なっている.
グラフ可視化手法を対象とした先行研究~\cite{saito},~\cite{YST},~\cite{BETT},~\cite{C}では, 
２次元の可視化結果を生成し, ネットワーク分析の有効性を実証している.
本研究の可視化手法では, ~\cite{saito}を採用するため,ここで各手法を紹介する. 
しかし, 本研究の提案法では, あくまでも「ばねモデル法」を採用しており,
他の可視化手法に関しては, 6章のネットワーク可視化評価法で取りあげるのみに留まっている.
すなわち, ばねモデル法以外の手法に関する理論の詳細は本研究の本筋ではないため, 詳細を省く.

\section{スペクトラル法}

隣接するノード対（$a_{i,j}=1$）の座標が近くなるように可視化する手法である.  

$x_i$,$x_j$ を任意のノードの座標ベクトルとし, $\|{\bf x}\|^2 = 1$ の制約のもとで, 次式を最小化する座標ベクトルを求める.  
\begin{eqnarray}
{\cal S}_1({\bf x}) & = & \frac{1}{2} \sum_{i=1}^{N} \sum_{j=1}^{N} a_{i,j} (x_i - x_j)^2. 
\label{spectral01}
\end{eqnarray}

\section{多次元尺度法}

高次元空間での複数観点の類似性を考慮し, 低次元空間へ可視化する手法である.

$H$ を, 任意の $N$ 次元ベクトルに対し, その要素の総和が $0$ になるよう平行移動させる行列としたとき,
$\|{\bf x}\|^2 = 1$ の制約のもとで, 次式を最大化する座標ベクトルを求める. 
\begin{eqnarray}
{\cal M}({\bf x}; {\bf G}) =  {\bf x}^T (-\frac{1}{2}{\bf H} {\bf G} {\bf H}) {\bf x}.
\label{MDS1}
\end{eqnarray}
ノード系列 $(v_{i_0}, v_{i_1}, \cdots, v_{i_k})$ において, 
隣り合う各ノード対 $(v_{i_{j-1}}, v_{i_j})$ にリンクがあるとき,
すなわち($a_{i_{j-1},i_j} = 1$)である場合, 
ノード $v_{i_0}$ から $v_{i_k}$ へのパスが存在すると言い, 
そのパスのリンク数 $k$ をパス長と呼ぶ. 
ここで, ノード $v_i$ と $v_j$ 間の距離（非類似度）を最短パス長 $g_{i,j}$ で定義すれば, 
全ノード集合を対象とした $N \times N$ の非類似度行列 ${\bf G}$ を考えることができる. 

\section{クロスエントロピー法} 

隣接ノードを非隣接ノードより近くに配置するために,
類似度間の交差エントロピーを用いて近似表現することで可視化する手法である.
次式を最小化する座標ベクトルを求める. 
\begin{eqnarray}
{\cal C}({\bf x}) & = & \sum_{i=1}^{N-1} \sum_{j=i+1}^{N} c(x_i, x_j), \label{cross} \\
c(x_i, x_j) & = & - a_{i,j} \log \rho(x_i, x_j) \label{cross2} \\
            &   & - (1-a_{i,j})\log(1-\rho(x_i, x_j)), \nonumber\\
\rho(x_i, x_j) & = & \exp(-\frac{1}{2} \|x_i - x_j\|^2). \label{neibor}
\end{eqnarray}
$\rho(x_i, x_j)$ は座標点間の類似度を定義し, 
この類似度と隣接行列要素との交差エントロピーを $c(x_i, x_j)$ で計算することで, 
任意のノード対での交差エントロピーの総和が最小化される. 

\section{ばねモデル法}

全てのエッジを「ばね」とみなし、「ばね」が伸縮しない自然長をエッジに設定することで、
ノードの重なりや不自然なエッジ長を抑制する可視化法である\cite{kamada89}.

$\beta_{i,j}$ をノード i, j を繋ぐエッジに対する「ばね定数」の初期値,
ノード $i, j$ 間の最短パス長を $g_{i,j}$ と定義し,
次式を最小化する座標ベクトルを求める. 
\begin{eqnarray}
{\cal K}({\bf x}) & = & \sum_{i=1}^{N-1} \sum_{j=i+1}^{N} \kappa(x_i, x_j), \label{spring} \\
\kappa(x_i, x_j) & = & \beta_{i,j}(g_{i,j} - \|x_i - x_j\|)^2. 
\label{spring2}
\end{eqnarray}
「ばね定数」 $\beta_{i,j}$ は $\beta_{i,j}= \frac{1}{2g_{i,j}^2}$ に設定され, 
$g_{i,j}$ が大きいほど, 自乗誤差の影響を小さくする働きを持つ. 

直感的には, 全エッジに「ばね」が装着されているとみなし, 
「ばね」が伸び縮みしない自然長を, ノード間の最短パス長 $g_{i,j}$ に設定すれば, 
「ばね」は自然長に近くなるように延び縮みするため, 不自然なエッジ長がなくなり,
結果的にエッジ両端のノードの重なりも抑制されるイメージである.

\newpage

\chapter{提案分析法}
提案分析法では,  クエリ文書と各検索結果文書から算出した, それぞれのコサイン類似度ベクトルを用い, 
(クエリ文書, 任意の検索結果文書)間の$D$-次元類似貢献度ベクトルを, 検索結果文書数だけ構築する.  
次に, 各検索結果文書をカテゴリー分けするために, 
それぞれが持つ$D$-次元類似貢献度ベクトルの各要素に着目し, 
その要素の最大値を用いて, $K$-最近傍グラフを構成する.
最後に, 構成したグラフを「ばねモデル法」で可視化するとともに, 
検索結果文書であることを表す各ノードを, 所属カテゴリーごとに配色し,
カテゴリー名に相当する「説明語」なるものを付与する.
以下では, $D$-次元類似貢献度ベクトル構築法と, $K$-最近傍グラフ構成法の詳細について述べる. 

\section{$D$-次元類似貢献度ベクトル構築法}

総数 $N$ の文書集合を ${\cal N} = \{1, \cdots, n, \cdots, N \}$ とし, 
語彙数 $D$ の出現単語集合を ${\cal D} = \{1, \cdots, d, \cdots, D \}$ とする. 
また, 各文書 $n$ とクエリ文書の $D$-次元特徴ベクトルを ${\bf x}_n$ と ${\bf q}$ とし,  
文書間でのコサイン類似度を求めるため, 特徴ベクトルはノルムが $1$ に正規化されているとする. 

まず, クエリ ${\bf q}$ に対し, 文書集合 ${\cal N}$ から $h$-最類似文書集合 ${\cal R}(h) \subset {\cal N}$ を求める. 
次に, 類似文書 $n \in {\cal R}(h)$ に対し, 特徴ベクトル ${\bf x}_n$ から単語 $d$ の類似貢献度を
\begin{eqnarray}
y_{n,d} & = & \frac{q_d \times x_{n,d}}{\sum_{d' \in {\cal D}} q_{d'} \times x_{n,d'}} 
\label{sc3-1}
\end{eqnarray}
で求め, $D$-次元類似貢献度ベクトル ${\bf y}_n$ を構築する. 
ここで, $y_{n,1} + \cdots + y_{n,D} = 1$ となる. 

\subsection{計算例}

総数 $6$ の文書集合を ${\cal N} = \{1, \cdots, n, \cdots, 6 \}$, 
語彙数 $D = 6$ の出現単語集合を ${\cal D} = \{車, 店, 鹿, 海, 山, 川\}$,
各文書の出現単語を, 
\begin{eqnarray}
doc_1 &=& (\text{車}, \text{車}, \text{車}, \text{店}, \text{店}, \text{鹿}, \text{海}, \text{山}) \\ 
doc_2 &=& (\text{車}, \text{店}, \text{店}, \text{店}, \text{鹿}, \text{鹿}, \text{海}, \text{川}) \\ 
doc_3 &=& (\text{車}, \text{店}, \text{店}, \text{鹿}, \text{鹿}, \text{鹿}, \text{山}, \text{川}) \\ 
doc_4 &=& (\text{車}, \text{店}, \text{店}, \text{海}, \text{海}, \text{海}, \text{山}, \text{川}) \\ 
doc_5 &=& (\text{車}, \text{鹿}, \text{海}, \text{山}, \text{山}, \text{山}, \text{川}, \text{川}) \\ 
doc_6 &=& (\text{店}, \text{鹿}, \text{海}, \text{山}, \text{山}, \text{川}, \text{川}, \text{川}) 
\label{sc3-2}
\end{eqnarray}
と定義し, $D$-次元類似貢献度ベクトル ${\bf y}_n$ の計算例を提示する.

まず, $D$-次元特徴ベクトル ${\bf x}_n$, ${\bf q}$をコサイン類似度を用いて
求めるために, 式~(\ref{sc3-3})から, 各文書の単語出現頻度を要素に持つ TF-IDF ベクトルを計算する.
なお, $doc_n$ は任意の文書を示す.

\begin{eqnarray}
\text{IDF}(d) = \log\left(\frac{N}{1 + \text{DF}(d)}\right)\\
\text{TF-IDF}(d, doc_n) = \text{TF}(d, doc_n) \cdot \text{IDF}(d)
\label{sc3-3}
\end{eqnarray}

今回の例では, 計算処理の把握を容易にするために,
IDF値を, $A = \log\left(\frac{6}{5}\right)$と定義する.
すると, 各文書の TF-IDFのベクトル値は以下となる.

\begin{eqnarray}
\text{TF-IDF}(doc_1) &=& (3, 2, 1, 1, 1, 0)A \\ 
\text{TF-IDF}(doc_2) &=& (1, 3, 2, 1, 0, 1)A \\ 
\text{TF-IDF}(doc_3) &=& (1, 2, 3, 0, 1, 1)A \\ 
\text{TF-IDF}(doc_4) &=& (1, 2, 0, 3, 1, 1)A \\ 
\text{TF-IDF}(doc_5) &=& (1, 0, 1, 1, 3, 2)A \\ 
\text{TF-IDF}(doc_6) &=& (0, 1, 1, 1, 2, 3)A
\label{sc3-4}
\end{eqnarray}

次に, 各 TF-IDF ベクトルの内積を下記で求める.

\begin{eqnarray}
\|\text{TF-IDF}(\mathbf{doc_1})\|^2 &=& (3^2+2^2+1^2+1^2+1^2+0^2)A^2 = 16A^2\\
\|\text{TF-IDF}(\mathbf{doc_1})\| &=& 4A\\
\|\text{TF-IDF}(\mathbf{doc_2})\|^2 &=& (1^2+3^2+2^2+1^2+0^2+1^2)A^2 = 16A^2\\
\|\text{TF-IDF}(\mathbf{doc_2})\| &=& 4A\\
\|\text{TF-IDF}(\mathbf{doc_3})\|^2 &=& (1^2+2^2+3^2+0^2+1^2+1^2)A^2 = 16A^2\\
\|\text{TF-IDF}(\mathbf{doc_3})\| &=& 4A\\
\|\text{TF-IDF}(\mathbf{doc_4})\|^2 &=& (1^2+2^2+0^2+3^2+1^2+1^2)A^2 = 16A^2\\
\|\text{TF-IDF}(\mathbf{doc_4})\| &=& 4A\\
\|\text{TF-IDF}(\mathbf{doc_5})\|^2 &=& (1^2+0^2+1^2+1^2+3^2+2^2)A^2 = 16A^2\\
\|\text{TF-IDF}(\mathbf{doc_5})\| &=& 4A\\
\|\text{TF-IDF}(\mathbf{doc_5})\|^2 &=& (0^2+1^2+1^2+1^2+2^2+3^2)A^2 = 16A^2\\
\|\text{TF-IDF}(\mathbf{doc_5})\| &=& 4A
\label{sc3-5}
\end{eqnarray}

これら TF-IDF ベクトルの内積値を式(\ref{sc3-6})に用い,  
(クエリ文書$q$, 検索結果文書$r$)のコサイン類似度ベクトル $x_n$ を求める.

\begin{eqnarray}
\text{cos}(\mathbf{q}, \mathbf{r}) = \frac{\mathbf{q} \cdot \mathbf{r}}{\|\mathbf{q}\| \|\mathbf{r}\|}
= \frac{\sum_{i=1}^D q_i r_i}{\sqrt{\sum_{i=1}^D q_i^2} \sqrt{\sum_{i=1}^D r_i^2}}
\label{sc3-6}
\end{eqnarray}

例えば, クエリ文書を $q = doc_1$, 検索結果文書を $r = doc_n$ としたとき, 
それぞれのコサイン類似度は以下となる.

\begin{eqnarray}
x_1 &=& \text{cos}(doc_1, doc_1) = \frac{(3, 2, 1, 1, 1, 0)A \cdot (3, 2, 1, 1, 1, 0)A}{4A \cdot 4A} = \frac{16}{16}\\
x_2 &=& \text{cos}(doc_1, doc_2) = \frac{(3, 2, 1, 1, 1, 0)A \cdot (1, 3, 2, 1, 0, 1)A}{4A \cdot 4A} = \frac{12}{16}\\
x_3 &=& \text{cos}(doc_1, doc_3) = \frac{(3, 2, 1, 1, 1, 0)A \cdot (1, 2, 3, 0, 1, 1)A}{4A \cdot 4A} = \frac{11}{16}\\
x_4 &=& \text{cos}(doc_1, doc_4) = \frac{(3, 2, 1, 1, 1, 0)A \cdot (1, 2, 0, 3, 1, 1)A}{4A \cdot 4A} = \frac{11}{16}\\
x_5 &=& \text{cos}(doc_1, doc_5) = \frac{(3, 2, 1, 1, 1, 0)A \cdot (1, 0, 1, 1, 3, 2)A}{4A \cdot 4A} = \frac{8}{16}\\
x_6 &=& \text{cos}(doc_1, doc_6) = \frac{(3, 2, 1, 1, 1, 0)A \cdot (0, 1, 1, 1, 2, 3)A}{4A \cdot 4A} = \frac{6}{16}
\label{sc3-7}
\end{eqnarray}

最後に, これらコサイン類似度を用いて, 式~(\ref{sc3-1})から, 
$x_n$の類似度貢献度ベクトル $y_n$ を算出する.

\begin{eqnarray}
y_1 &=& \frac{
\left( 
\frac{3}{4} \cdot \frac{3}{4}, 
\frac{2}{4} \cdot \frac{2}{4}, 
\frac{1}{4} \cdot \frac{1}{4}, 
\frac{1}{4} \cdot \frac{1}{4}, 
\frac{1}{4} \cdot \frac{1}{4}, 
\frac{0}{4} \cdot \frac{0}{4} 
\right)
}{\frac{16}{16}} 
= \left( \frac{9}{16}, \frac{4}{16}, \frac{1}{16}, \frac{1}{16}, \frac{1}{16}, \frac{0}{16} \right)\\
y_2 &=& \frac{
\left( 
\frac{3}{4} \cdot \frac{1}{4}, 
\frac{2}{4} \cdot \frac{3}{4}, 
\frac{1}{4} \cdot \frac{2}{4}, 
\frac{1}{4} \cdot \frac{1}{4}, 
\frac{1}{4} \cdot \frac{0}{4}, 
\frac{0}{4} \cdot \frac{1}{4} 
\right)
}{\frac{12}{16}} 
= \left( \frac{3}{12}, \frac{6}{12}, \frac{2}{12}, \frac{1}{12}, \frac{1}{12}, \frac{0}{12} \right)\\
y_3 &=& \frac{
\left( 
\frac{3}{4} \cdot \frac{1}{4}, 
\frac{2}{4} \cdot \frac{2}{4}, 
\frac{1}{4} \cdot \frac{3}{4}, 
\frac{1}{4} \cdot \frac{0}{4}, 
\frac{1}{4} \cdot \frac{1}{4}, 
\frac{0}{4} \cdot \frac{1}{4} 
\right)
}{\frac{11}{16}} 
= \left( \frac{3}{11}, \frac{4}{11}, \frac{3}{11}, \frac{0}{11}, \frac{1}{11}, \frac{0}{11} \right)\\
y_4 &=& \frac{
\left( 
\frac{3}{4} \cdot \frac{1}{4}, 
\frac{2}{4} \cdot \frac{2}{4}, 
\frac{1}{4} \cdot \frac{0}{4}, 
\frac{1}{4} \cdot \frac{3}{4}, 
\frac{1}{4} \cdot \frac{1}{4}, 
\frac{0}{4} \cdot \frac{1}{4} 
\right)
}{\frac{11}{16}} 
= \left( \frac{3}{11}, \frac{4}{11}, \frac{0}{11}, \frac{3}{11}, \frac{1}{11}, \frac{0}{11} \right)\\
y_5 &=& \frac{
\left( 
\frac{3}{4} \cdot \frac{1}{4}, 
\frac{2}{4} \cdot \frac{0}{4}, 
\frac{1}{4} \cdot \frac{1}{4}, 
\frac{1}{4} \cdot \frac{1}{4}, 
\frac{1}{4} \cdot \frac{3}{4}, 
\frac{0}{4} \cdot \frac{2}{4} 
\right)
}{\frac{8}{16}} 
= \left( \frac{3}{8}, \frac{0}{8}, \frac{1}{8}, \frac{1}{8}, \frac{3}{8}, \frac{0}{8} \right)\\
y_6 &=& \frac{
\left( 
\frac{3}{4} \cdot \frac{0}{4}, 
\frac{2}{4} \cdot \frac{1}{4}, 
\frac{1}{4} \cdot \frac{1}{4}, 
\frac{1}{4} \cdot \frac{1}{4}, 
\frac{1}{4} \cdot \frac{2}{4}, 
\frac{0}{4} \cdot \frac{3}{4} 
\right)
}{\frac{6}{16}} 
= \left( \frac{0}{6}, \frac{2}{6}, \frac{2}{6}, \frac{1}{6}, \frac{1}{6}, \frac{2}{6} \right)
\label{sc3-8}
\end{eqnarray}

類似度貢献度ベクトル $y_n$ における各要素の最大値を抽出すると, 
$doc_n$ の説明語 $d_n$ が得られる.\\
$d_1$ : 車\\
$d_2$ : 店\\
$d_3$ : 店\\
$d_4$ : 店\\
$d_5$ : 車\\
$d_6$ : 店

\section{グラフ構成法}

類似貢献度ベクトル ${\bf y}_n$ 間の距離に基づき, $h$-最類似文書集合 ${\cal R}(h) \subset {\cal N}$ をノード集合として, 
$K$-最近傍 (NN: Nearest Neighnour) グラフ $G({\cal R}(h), {\cal E})$ を構築する. 
当手法は, 与えられたデータ点集合に基づき, 各点がその近傍に位置する$K$個の点とエッジで結ばれるグラフを構成するものである.
下記手順で$K$-最近傍グラフを構築した.
\begin{enumerate}
    \item クエリ文書${\bf q}$に対して最も類似貢献度が高い$k$個の$y_n$をエッジで接続.
    \item $h$-最類似文書集合の各点でも同様の処理を$h$回実行.
\end{enumerate}
最後に, 単語 $d$ で類似貢献度が最大となるノード(文書)集合を
\begin{eqnarray}
{\cal C}(d) & = & \{ n \in {\cal R}(h)~|~d = \argmax_{d' \in {\cal D}} (y_{n,d'}) \} 
\label{sc2}
\end{eqnarray}
として求め, そこに属すノード群を同色にし, 単語 $d$ を説明語として表示可能にした 
$k$-最近傍 $G({\cal R}(h), {\cal E})$ の可視化結果を出力する. 

また, 類似貢献度ベクトル間に通常のユークリッド距離を導入し,
ラベル集合 ${\cal E}$ に対して, 最小全域木(MST: Minimum Spanning Tree)を構築すれば, 
グラフ $G_{\mbox{MST}} = ({\cal E}, {\cal F}_{\mbox{MST}})$ を得ることができる. 
ここで, ${\cal E}$ はラベルを要素としたノード集合に対応し,
${\cal F}_{\mbox{MST}}$ は MST を構成するエッジ集合を表す. 

\section{グラフ可視化法}

グラフ可視化には多様な方法~\cite{saito}を適用できるが, 
本研究では, 重なるノード数が最小になるという理由で, ばねモデル法~\cite{kamada89}を採用した.
また, そもそもグラフを用いずに可視化するという選択も存在するが, 
その場合と比較すると, グラフ可視化では類似性の高いノード (文書) のペアにエッジが接続されるので, 
2次元空間への埋め込みにおいて偶発的に近くに配置されたケースとの区別が容易にできる.
なお, 類似性の度合いに比例するようにエッジ長を調整することで, より適切な可視化結果が得られことが期待できるが, 
重なるノードが多数出現する, 不自然なエッジ長が出現する等の問題が発生するため,
本研究では重み無しグラフの範囲を扱い, 重み付きグラフ可視化への拡張は今後の課題とする. 
これら問題については6章にて詳細を明記する. 

\chapter{評価実験}
本章では, まず, 実験データについて説明する.
次に, 提案法での可視化結果を示す.
最後に, 提案法から生成した可視化グラフが視覚的に把握可能であることを示す.

\section{実験データ}

本実験では，文書数 $N = 7,367$ の livedoor ニュースコーパスを文書集合 ${\cal N}$ とし, 
MeCab~\cite{kudo04}で出力される形態素をすべて単語として扱い, 語彙数は $D = 89,993$ とした.
なお, グラフ可視化には, ばねモデル法~\cite{kamada89}を採用し, 
これらの設定を元に, それぞれ文書内容や容量が異なる検索クエリ文書1,2,3を用いて生成したグラフ可視化結果を説明する.
また, 検索クエリ文書1,2,3の全内容については以下に示す.

\newpage
\subsection{検索クエリ文書1}

ゆるキャラ、ネクストジェネレーション
独女のココロを優しく癒すゆるキャラたち。ドアラ、ひこにゃん、せんとくんなどなど、そのクオリティの如何に関わらず、なぜか「大ブレイク」してしまうところもゆるキャラの魅力。一度火がつけばその人気は一気に全国区となり、関連商品は飛ぶように売れまくる。

そのゆるキャラ界で新たなスターとなりそうな1匹（？）こそ、ダイハツ自動車「ムーヴコンテ」のCMに登場している「カクカクシカジカ」だろう。コミカルな声にあわせ「カクカクシカジカ」と軽妙な動きを見せるシカ——、といえば、一度は目にしたこともあるはず。

初登場は、新車ムーヴコンテが発表された08年8月のCMで、「このカクカクしたクルマ、いいところがシカジカ。ゆえにカクカクシカジカ」と、クルマの説明なのか、自己紹介なのかすらも不明なセリフで登場。以降、自分がシカであるくせに、干支の牛や虎になるために気ぐるみを着てみたり、ムーヴコンテに搭載されている新技術の「CVT」を「ちょっと ヴぃっくりな 低燃費」と解説するなど、ちょっとひねくれた辛口評論家という設定を無視するかのごとく、キャラが独走を続けている。

そのゆるキャラのネクストジェネレーションであるカクカクシカジカのオリジナルグッズがもらえるキャンペーン「カクカクトクトクフェア」が、6月5日（土）、6日（日）の両日で開催される。期間中にダイハツのお店に来店すると、抽選で「カキ氷器」「折りたたみ傘」「ペットッボトルホルダー」がもれなく当たるという。

フェア告知用のCMも製作されているようで、今回、カクカクシカジカが扮するのはサッカーの審判。ワールドカップというタイミングに合わせての演出のようだが、今作も適当なゆるさを楽しませてくれるはず。CMは本日よりOA。

「かわいらしく、スタイリッシュな軽自動車」というブランドイメージを持つダイハツが、さらに女性からの支持率を上がりそうなゆるキャラの躍進。この週末には普段はなかなか行く機会がないディーラーにゆるキャラ目当てで行ってみるのもいいかもしれない。

・カクカクトクトクフェア

・奔放な振る舞いを見られる過去のCMアーカイブ【コチラ】

■「カクカクトクトクフェア」
2010年6月5日（土）〜6月6日（日）
全国のダイハツのお店
※一部のお店では実施していない場合がございます。

■景品
1等 カクカクシカジカ　カキ氷器
2等 カサカサシカジカ　折りたたみ傘
3等 ゴクゴクシカジカ　ペットボトルホルダー

\newpage
\subsection{検索クエリ文書２}

【Office Scope vol.09】ゴー・アブロード！　世界一周できるオフィス（フォートラベル株式会社）
　ワークスペースのデザイニングに優れた企業のオフィスを訪問し、そこで働く人々の声から、その”働き心地”をリポートする連載企画『Office Scope』。第9回の今回は、カカクコムグループで「旅行のクチコミサイト フォートラベル」を運営する、フォートラベル株式会社へ訪問。旅行先選びから思い出の共有まで———旅行に関わる全ての満足を、独自の視点と多彩なサービスで実現する同社。そんな同社のオフィス内コミュニケーションはいかに？
 
エントランスは空港のイメージで「これから旅立つぞ」というワクワク感が。
創業当初、オフィスはファストフード店だった!?

　「旅行のクチコミサイト フォートラベル」の運営を行っている、フォートラベル株式会社。2003年の創業より、旅行者のクチコミをもとに充実した旅行情報を提供するサービスを行っている。2005年より株式会社カカクコムのグループ会社に加入。数回のオフィス移転を経て現在のオフィスへ入ったが、設立当初はオフィスを持たずにファストフード店などで地道に作業を行っていた時期もあったという。さらに、オフィスを持ってからもなかなか思い通りにはいかなかったようで・・・。
「ここに入る以前のオフィスは、ミーティングスペースもないような小さな場所だったので、打ち合わせは喫茶店などで行うこともしばしばありました。狭かったので、席を立つと後ろの席の人にぶつかるようなことも・・・」（フォートラベル株式会社広報）
現在のオフィスではミーティングルームが5部屋あるほか、執務室には簡単な打ち合わせを行うことができるようなスペースが設けられている。わざわざ外へ出なくとも話し合いができるので、仕事の効率は上がったようだ。また、社員用のデスクは区切りのない一枚板になっているため、隣の席との境界がなくフラット。こちらもコミュニケーションを円滑にする役割も担っているそうだ。

執務室に設けられたミーティング用のデスクの上には、会社のロゴになっている“4”のマークが。

白いデスクはオーダーメイドだそう。赤いキャビネットが特徴的。

世界一周気分を味わえるミーティングルーム
　5つあるミーティングルームは世界の都市をイメージして設えられ、中国の「上海」、ギリシャの「サントリーニ島」、フィンランドの「ヘルシンキ」、インドネシアの「バリ島」、日本の「嵯峨野」と、それぞれテーマ性の強い空間になっている。念願のミーティングルームとだけあって、社内でアイデアを出し合って決めたのだという。
「飛行機や新幹線のイメージという案も出たのですが、話し合った結果、各国の特徴的な要素を落とし込むということに。デザイン会社の方と相談しながら、部屋のコンセプトを決めたり、それに合った壁紙などを選んだりしていきました」
　それぞれの部屋にはホワイトボードが備え付けてあるが、カーテンやブラインドで隠れるように工夫がなされている。日本にいながら、ワールドワイドな非日常感を味わうことができる。

上海はハイ・グレードホテルの一室のようなデザイン。

こちらはギリシャのサントリーニ島をイメージ。白亜に青のコントラストが美しい。

フィンランドのヘルシンキを体現したデザインは、映画「かもめ食堂」に着想を得たのだとか。

ソファでつい長居したくなりそうな、インドネシアのバリ島をイメージした一室。

日本からは、京都の嵯峨野。茶だんすやレトロな時計が趣深い。

今後は、オフィスをもっと賑やかに・・・ 
　オフィス内にはミーティングルームのほか、リフレッシュルームやライブラリ（資料閲覧スペース）も設けられている。ちなみにライブラリは、メキシコにある「ルイス・バラガン邸」という有名な建築物を模したカラフルな部屋になっている。こちらには「地球の歩き方」や、雑誌版「4travel」といった紙媒体の旅行資料が詰まっている。
　ここまで見る限りほぼ完成形に近いような気がするが、オフィスについて今後どうしていきたいかというと・・・？
「弊社には“旅行休暇”といって、土日と合わせると最大5日間休みをとれる制度があるんです。1日あたり1万円の補助金が出るので、それを活用して社員は旅行を楽しんでいます。社員が旅行先で購入した現地の民芸品などをミーティングルームに飾っていって、ちょっとずつ賑やかにしていければと思います」
　現地色の強い小物が増えれば、さらに各国の雰囲気が出るに違いない。

「ルイス・バラガン・モルフィン」というメキシコ人の建築家がデザインした建築物にインスパイアーされたライブラリ。

バリ島のミーティングルームに設えられたカエルの置物からは水が・・・。

こちらの小物は実際に旅行した社員が購入したお土産。今後もっと増えていくかも？

リフレッシュルームには、社員からの旅行のお土産が並ぶことも多いとか。

「旅行のクチコミサイト フォートラベル」の活用方法も聞いてみました 
　旅行者のために有益な情報を提供している「旅行のクチコミサイト フォートラベル」。その上手な活用方法を伺った。
「バックパッカーからハイエンドな旅行を好む方まで、あらゆるユーザーに登録していただいています。旅行の情報について知りたいことがあるときは『教えて！トラベラー』というQ＆A掲示板で質問してみると、旅行に詳しいユーザーから回答を受けることができます。また旅行先を指定すると、会員登録した自分の管理画面にその土地の最新情報が表示されるなどの機能もあります」
今後は、さらにクチコミの投稿を充実させることを目標にしているという。
「旅行には価値観が出るから、さまざまなユーザーの方からクチコミをいただければ探し物をしている人にとって便利です。情報交換の場としてもぜひ活用してほしいですね」
　旅行を考えている人はもちろん、旅行欲を高めたい人、現実からちょっと目を逸らしたい人にも（!?）おすすめのサイトだ。

オープンスペースはチェコのプラハをイメージ。ユーザーを招いたこともあるという。

 ミーティングルームの外も街道のような凝ったつくりになっている。

■関連リンク
　・旅行のクチコミサイト フォートラベル
　・Office Scope過去記事一覧

livedoor求人・転職
livedoor求人・転職は、あなたがどんな求人情報を探しているのかを瞬時に判断しておすすめするコンテンツ。「東京都内で年収700万以上の営業職で未経験OK」などあなたのこだわりにマッチする転職活動を完全サポート。
http://type.jp/s/livedoor/

\newpage
\subsection{検索クエリ文書３}

ベイビースキンを手に入れる！　HELENA RUBINSTEIN×THE WESTIN TOKYO 初肌〜うぶはだ〜スパ
　夏が終わり、そろそろ「紫外線」での肌荒れをケアして、冬の乾燥に備えたい女子も多いはず。15周年を迎えるウェスティンホテル東京（目黒区三田1-4-1）は、世界初の特別スパトリートメントとレディース宿泊プランHELENA RUBINSTEIN×THE WESTIN TOKYO 初肌〜うぶはだ〜スパを10月31日まで提供している。

　プロディジー リプラスティ ピール シリーズを贅沢に使った「ヘレナ ルビンスタイン」特別スパトリートメントでは、専用のラグジュアリースイート（通常1 室/8,400）で日々の紫外線やストレスなどで代謝の鈍った肌のターンオーバーを正常に戻し、目覚めるような初肌（うぶはだ）を体感できる。

　また、高級感溢れる特別アメニティ、エグゼクティブクラブ ラウンジで朝食やティータイムを過ごす事ができるのも嬉しい。女性専用フロアでのエグゼクティブステイで、ラグジュアリー気分に酔いしれながら、心も体もリニューアルを体感できる、忙しい女子にとって必見のプランである。

・詳細
期間：2009年9月4日（金）〜10月31日（土）
特典： プラン限定アメニティプレゼント、エグゼクティブ クラブ ラウンジの利用、レイトチェックアウト14:00（通常12:00）

・お問合せ先
03-5423-7770（9：00〜19：00）宿泊予約課

■関連リンク
・HELENA RUBINSTEIN×THE WESTIN TOKYO 初肌〜うぶはだ〜スパ

\newpage

\section{提案法での可視化結果}

まず図~\ref{fig:3-1}に，提案法での可視化結果の例を示す. 
比較的大きな赤丸は検索クエリ文書である. 
類似文書数は $h = 100$ に設定し, $h$-最類似文書集合 ${\cal R}(h)$ を求めた.  
ノード配色は, ${\cal R}(h)$ を類似度順で降順に並べ, 
同じ説明語となるノード群は同色となるよう, CMYKカラーシステム24色をサイクリックに割り当てた. 

図~\ref{fig:3-1}(a)に示す, 検索クエリ文書１を用いた可視化結果からは,  
CM, 店, クルマを説明語とし, 比較的密結合する多数のノード群とともに,  
比較的大きな赤丸のクエリ近くに, 少数ながらダイハツ, ゆるキャラを説明語とするノード群も見て取れる. 
その他の単一説明語としては, 自動車, キャンペーンなどが出現する. 
また, 図~\ref{fig:3-1}(b)に示す, 検索クエリ文書２を用いた可視化結果では,  
旅行, オフィス, 転職を説明語とした, 密結合する多数のノード群が構成されており,
その他の説明語としては社員が出現する. 
最後に, 図~\ref{fig:3-1}(c)に示す, 検索クエリ文書３を用いた可視化結果からは,  
肌, 紫外線, プランを説明語とした密結合ノード群が構成されている. 
これら結果より, ある観点で類似する文書数がどの程度存在するか視覚的に把握可能と言える. 

\begin{figure}
   \begin{center}
    \includegraphics[width=0.65\hsize]{figs/fig2a.png}
    \mbox{(a)検索クエリ文書１を用いて生成した可視化結果}
   \end{center}
   \begin{center}
    \includegraphics[width=0.65\hsize]{figs/ku/fig3-1/fig1b.png} 
    \mbox{(b)検索クエリ文書２を用いて生成した可視化結果}
   \end{center}
   \begin{center}
    \includegraphics[width=0.65\hsize]{figs/ku/fig3-1/fig1c.png}
    \mbox{(c)検索クエリ文書３を用いて生成した可視化結果}
   \end{center}
 \caption{提案法を用いて生成した可視化グラフ}
 \label{fig:3-1}
\end{figure}

\section{検索文書数による視覚的効果確認}

次に, 検索文書数を変更して, グラフの視認性への影響を確認する.
図~\ref{fig:3-2}に, 検索文書数を $50, 100, 200$ に設定した可視化グラフを示す.
検索文書数が50の(a)は, 同グループごとの文書集合は明確に分けられて可視化されているが, 
文書数が少なすぎることで蜜結合部分が視覚的に分かりずらく, 部分構造の把握に難点がある.
検索文書数が200の(b)は, 一部のノード結合が異なる説明語(配色)となり, 全体的な類似構造把握の観点でも限界が見受けられる.
このようになってしまう原因として, 母集団となる7367個の文書集合に類似している文書が200件も存在しないにも関わらず,
無理矢理200件検索したことで, 関連性の薄い文書まで可視化されてしまったことが挙げられる.
これら結果より, 当実験データを用いた可視化グラフでは, 今後は検索文書数を $h = 100$ としたものを対象とする.

\begin{figure}
   \begin{center}
    \includegraphics[width=0.75\hsize]{figs/ku/fig3-2/fig1b.png}
    \mbox{(a)$h=50$ に設定した可視化結果}
   \end{center}
   \begin{center}
    \includegraphics[width=0.75\hsize]{figs/ku/fig3-2/fig1c.png} 
    \mbox{(b)$h=200$ に設定した可視化結果}
   \end{center}
   \begin{center}
    \includegraphics[width=0.75\hsize]{figs/ku/fig3-2/fig1a.png}
    \mbox{(c)$h=100$ に設定した可視化結果}
   \end{center}
 \caption{検索文書数の変更に伴う可視化結果の例}
 \label{fig:3-2}
\end{figure}

\chapter{他手法との比較}
本章では, まず, 比較法から生成した可視化グラフは視覚的な把握が困難であることを確認する. 
次に, 比較法と提案法それぞれで生成した可視化グラフを示し, 
提案法から生成した可視化グラフの方が視覚的把握が容易であることを示す.
最後に, 上記の定性評価を定量的にも評価することで, 提案法の有効性を示す. 

\section{比較法の定義}
比較評価のため, 式~(\ref{sc3-1})でクエリとの類似貢献度を考慮せず, 
類似文書 $n \in {\cal R}(h)$ それぞれの類似度として, 
\begin{eqnarray}
z_{n,d} = x_{n,d}^2
\label{sc4-1}
\end{eqnarray}
を求め, 
$D$-次元ベクトル ${\bf z}_n$ を構築し, 説明語を $\argmax (z_{n,d'})$で求める方法を比較法と呼ぶ. 
なお, 先行研究~\cite{oda04}では, クエリを考慮した類似貢献度ベクトルを用いず, 
検索結果文書間の類似度を求める点で比較法の範疇と見なせる. 
また, 提案法において, グラフ構築に最小全域木(MST: Minimum Spanning Tree)を用いる方法をMST法と呼び比較評価する.

\section{構成グラフの定性評価}

図~\ref{fig4-1}(b)に示す比較法での可視化結果では, 
殆どのノード結合は異なる説明語 (配色) となり, 全体的な類似構造把握の点などでも限界が見受けられる. 
図~\ref{fig4-2}, ~\ref{fig4-3}に関しても同様の結果となった.

\begin{figure}[H]
   \begin{center}
    \includegraphics[width=\hsize]{figs/fig2a.png}
    \mbox{(a) クエリ１から生成した提案法での可視化結果}
   \end{center}
   \begin{center}
    \includegraphics[width=\hsize]{figs/fig2b.png} 
    \mbox{(b) クエリ１から生成した比較法での可視化結果}
   \end{center}
 \caption{クエリ１から生成した提案法と比較法の可視化結果}
 \label{fig4-1}
\end{figure}

\begin{figure}[H]
   \begin{center}
    \includegraphics[width=\hsize]{figs/ku/fig3-1/fig1b.png}
    \mbox{(a) クエリ２から生成した提案法での可視化結果}
   \end{center}
   \begin{center}
    \includegraphics[width=\hsize]{figs/ku/fig4-2/fig1b.png} 
    \mbox{(b) クエリ２から生成した比較法での可視化結果}
   \end{center}
 \caption{クエリ２から生成した提案法と比較法の可視化結果}
 \label{fig4-2}
\end{figure}

\begin{figure}[H]
   \begin{center}
    \includegraphics[width=\hsize]{figs/ku/fig3-1/fig1c.png}
    \mbox{(a) クエリ３から生成した提案法での可視化結果}
   \end{center}
   \begin{center}
    \includegraphics[width=\hsize]{figs/ku/fig4-3/fig1c.png} 
    \mbox{(b) クエリ３から生成した比較法での可視化結果}
   \end{center}
 \caption{クエリ３から生成した提案法と比較法の可視化結果}
 \label{fig4-3}
\end{figure}

\newpage
一方, 図~\ref{fig4-4}に示すMST法での可視化結果と比較すれば, 
ノード配色は提案法と同一であるが, 密結合する部分構造把握の点などで, 提案法に望ましい性質が見て取れる.

\begin{figure}[H]
   \begin{center}
    \includegraphics[width=\hsize]{figs/fig2a.png}
    \mbox{(a) クエリ１から生成した提案法での可視化結果}
   \end{center}
   \begin{center}
    \includegraphics[width=\hsize]{figs/fig2c.png}
    \mbox{(b) クエリ１から生成したMST法での可視化結果}
   \end{center}
 \caption{クエリ１から生成した提案法とMST法の可視化結果}
 \label{fig4-4}
\end{figure}

\begin{figure}[H]
   \begin{center}
    \includegraphics[width=\hsize]{figs/ku/fig3-1/fig1b.png}
    \mbox{(a) クエリ２から生成した提案法での可視化結果}
   \end{center}
   \begin{center}
    \includegraphics[width=\hsize]{figs/ku/fig4-5/fig1b.png}
    \mbox{(b) クエリ２から生成したMST法での可視化結果}
   \end{center}
 \caption{クエリ２から生成した提案法とMST法の可視化結果}
 \label{fig4-5}
\end{figure}

\begin{figure}[H]
   \begin{center}
    \includegraphics[width=\hsize]{figs/ku/fig3-1/fig1c.png}
    \mbox{(a) クエリ３から生成した提案法での可視化結果}
   \end{center}
   \begin{center}
    \includegraphics[width=\hsize]{figs/ku/fig4-6/fig1c.png}
    \mbox{(b) クエリ３から生成したMST法での可視化結果}
   \end{center}
 \caption{クエリ３から生成した提案法とMST法の可視化結果}
 \label{fig4-6}
\end{figure}

\section{構成グラフの定量評価}

次に, 類似貢献度ベクトルを用いる効果を定量評価する. 
まず, $h$-最類似文書集合 ${\cal R}(h)$ における各単語 $d$ での類似貢献度の期待値を
\begin{eqnarray}
u_d = \sum_{n \in {\cal R}(h)} y_{n,d}/h
\label{sc4-2}
\end{eqnarray}
で求め, 
これら値を降順
\begin{eqnarray}
u_{d(r)} \geq u_{d(r+1)}
\label{sc4-3}
\end{eqnarray}
に並び替え, その累積値を
\begin{eqnarray}
v_s = \sum_{r=1}^s u_{d(r)}
\label{sc4-4}
\end{eqnarray}
で求める. よって, $v_D = 1$ となる. 

図~\ref{fig:sc}に，クエリ文書を文書集合 ${\cal D}$ のそれぞれに設定したときの累積値 $v_s$ の期待値を示す. 
ここで, p050, p100, p200 は上述した提案法で検索文書数を, $h = 50, 100, 200$ に設定した結果であり, 
c050, c100, c200 には, 類似貢献度期待値を特徴量での期待値
\begin{eqnarray}
u_d' = \sum_{n \in {\cal R}(h)} z_{n,d}/h
\label{sc4-5}
\end{eqnarray}
に置き換えた比較法において, $h = 50, 100, 200$ に設定した結果である. 
図~\ref{fig:sc}より，検索文書数 $h$ に依らず, 
提案法では上位 $r = 10$ 単語程度で累積貢献値が0.5程度になるのに対して, 
比較法では, そのために上位数百単語程度が必要になる. 
すなわち, 提案法では, 比較的少数の単語に類似貢献度が集中するので, 
少ない説明語数でも検索結果文書全体の大まかな内容が把握可能といえる.
従って, 妥当なノード配色や説明語選択の実現を期待できる. 

\begin{figure}[H]
  \begin{center}
  \includegraphics[width=\hsize]{figs/fig1.png}
  \end{center}
  \caption{貢献度期待値の累積分布}
  \label{fig:sc}
\end{figure}

\chapter{定量性を考慮した類似文書可視化分析法の評価}
我々は, 任意のノードをホバーすることで, 文書詳細を閲覧できる機能をもつ可視化グラフを想定し, 
与えられた文書クエリの類似検索結果文書集合を, グラフを用いて可視化する手法を提案している~\cite{ozaki1}.
一方で, グラフ可視化には他の手法~\cite{saito},~\cite{kamada89}を用いることも可能だが, 
可視化結果を定量評価する方法はあまり知られていない.
本研究では, 各種可視化法の有効性を, エッジ長の変動係数を用いて定量評価する方法を提案している~\cite{ozaki2}.
本章では，３章で示した提案手法を類似文書可視化分析法とし，
その手法を用いたグラフにおける可視化結果の有効性を, 
エッジ長の変動係数を用いて定量評価する方法を提案する.
本章では, ４つの手法でグラフを可視化し, 
定量評価とともに可視化例に基づく定性評価により, 定量評価法の妥当性を検証する. 

\section{評価手法}

任意のノードをホバーすることで, 文書詳細を閲覧できる機能をもつ可視化グラフを想定する.
まず, 提案法を述べる前に, 良い可視化グラフといえる基準を定義する.
本研究では下記項目を設定した.
\begin{enumerate}
    \item 隣接ノードが離れすぎずに描画される
    \item 重なっているノードが存在しない
\end{enumerate}
項目１は, 隣接ノード同士が一定の範囲内で描画されるか評価している.
非常に離れた位置に点在している場合, ホバーしにくいためである.
項目２は, 複数のノードが重なって描画されないか評価している.
重なるとホバーが困難なため, 最悪の場合, 文書詳細を閲覧できないノードが出現する.
従って, 上記項目をどの程度満たすか定量的に評価可能な手法を提案することが目的となる.
次に, 提案可視化評価法の手順を以下に述べる.
エッジ長の集合を ${\cal E} = \{1, \cdots, e, \cdots, E \}$ とし,
${\cal E}$ に対し, 平均 $\mu$, 標準偏差 $\sigma$ を求める.
次に, 変動係数を
\begin{equation}
  {cv} = \frac{\sigma}{\mu} \times 100
  \label{sc6-1}
\end{equation}
で導出する.
式(\ref{sc6-1})は, エッジ長のばらつき度合いを示す数値であり,
この値が小さい程, 各エッジの長さは全体的に同程度となる.
故に, 隣接ノードが一定の範囲内に描画される割合が高くなる.
逆に, 値が大きい程, 各エッジの長さに「ばらつき」が生じる.
故に, 短すぎる, 長すぎる等の不自然なエッジ長をもつ割合が大きくなり, 
結果的に隣接ノード同士の距離が不自然なペアが出現する.
従って, 変動係数が小さくなる程, 項目１を高水準で満たすグラフとなるため,
式(\ref{sc6-1})を用いることで, ある程度の定量評価が可能になる.

次に, 項目２に該当するケースも考慮する.
ここでは, 半径 $r$ の円として可視化する.
例えば, エッジの生成範囲が一方向に偏っている場合, 
ノードが重なる割合が高くなり, ホバーが困難なグラフが生成されてしまう.
従って, ノード集合を ${\cal N} = \{1, \cdots, n, \cdots, N \}$, 
ノード $n$ の座標を $(X_n, Y_n)$, $\forall n \neq m$ の条件で重なっているノードの個数を $on$とし,
\begin{equation}
  {on} = |\left\{ n, m \mid \sqrt{(X_n - X_m)^2 + (Y_n - Y_m)^2} \leq 2r \right\}|
  \label{sc6-2}
\end{equation}
を求め, 定量評価に用いる.
この値が小さくなる程, 項目２を高水準で満たすグラフとなる.

最後に, 可視化ネットワーク数を$L$, 可視化手法を$M$とした$L*M$個のデータを用い, 
X軸には評価項目１を式(\ref{sc6-1})で求めた値を,
Y軸には評価項目２を式(\ref{sc6-2})で求めた値を設定した散布図を生成する. 

\section{実験による評価}\label{sec:EXP}

提案法での評価項目１,２を用いて, 可視化グラフを定量評価可能か実験する.
可視化手法$M$には, スペクトラル法(se), 多次元尺度法(mds),
クロスエントロピー法(ce), ばねモデル法(kk)を採用した.
また, 定量評価する際の実験データとして, livedoor ニュースコーパスから文書検索を行い, 
ノード数を $100$ に設定した全7367件の可視化グラフを生成した. 
本稿では, 標本集団として抽出した$L=8$, $M=4$の範囲内で図~\ref{fig6-1}の散布図を用いて説明する.

\begin{figure}[tb]
  \begin{center}
  \includegraphics[width=\hsize]{figs/ipsj2025/fig1.png}
  \end{center}
  \caption{可視化グラフの評価値を示す散布図}
  \label{fig6-1}
\end{figure}

X軸には評価項目１を式(\ref{sc6-1})で求めた値を,
Y軸には評価項目２を式(\ref{sc6-2})で求めた値を設定した.
左下に近い位置にプロットされるサークルである程, 項目１,２を高水準で満たすグラフとなる.
また, 各サークルの数値はクエリ文書の $id$ を示しており, 同番のサークルは同様のクエリ文書を用いて可視化グラフを生成している.
先ほどの可視化例で示した図~\ref{fig:3-1}(a),(b)(c)は, クエリ文書$1,2,3$である.

図~\ref{fig6-2}より, ばねモデル法はどれも左下にプロットされており, 可視化グラフとして高水準である.
一方で, スペクトラル法は全体的に右上にプロットされているため低水準と判断できる.
多次元尺度法, クロスエントロピー法は全体的に図の中央に位置しており, 
それぞれを比較してみると, 項目１は多次元尺度法, 項目２はクロスエントロピー法に分があることが見て取れる.

次に, 実際に生成された可視化グラフを確認し, 先ほどの定量評価と結果が一致するか確認する. 
今回は, 各種可視化グラフの評価値が全体的に同程度となった,
すなわち, 目視での評価が困難である, クエリ文書$6$を対象として確認した.

クエリ文書$6$のグラフ可視化結果を図~\ref{fig6-2}に示す. 
スペクトラル法(a), 多次元尺度法(b), クロスエントロピー法(c)では,
離れている, 部分的に重なっている隣接ノードが存在するが,
ばねモデル法(d)では全く存在しないことが見て取れる. 
これら結果より, 定性評価でも, ばねモデル法が高水準であり, 定量評価と結果が一致した.

\begin{figure}[t]
   \begin{center}
    \includegraphics[width=0.6\hsize]{figs/ipsj2025/fig2a.png} 
    \mbox{(a)スペクトラル法}
   \end{center}
   \begin{center}
    \includegraphics[width=0.6\hsize]{figs/ipsj2025/fig2b.png}
    \mbox{(b) 多次元尺度法}
   \end{center}
   \begin{center}
    \includegraphics[width=0.6\hsize]{figs/ipsj2025/fig2c.png}
    \mbox{(c) クロスエントロピー法}
   \end{center}
   \begin{center}
    \includegraphics[width=0.6\hsize]{figs/ipsj2025/fig2d.png}
    \mbox{(d)ばねモデル法}
   \end{center}
 \caption{クエリ文書$6$の各種グラフ可視化結果}
 \label{fig6-2}
\end{figure}

\chapter{おわりに}
本論文では, 検索結果の各文書をノードと見なし, 
それぞれを類似貢献度ベクトルに変換し, 
k-最近傍グラフを構築した可視化において, 
検索結果文書がクエリに対し, どのような説明語で類似するかを付与するとともに, 
グループ分けして配色する新手法のグラフ可視化法を提案した. 

また, 可視化グラフの評価方法として, エッジ長の変動係数を定量評価する手法を採用することで, 
どの手法を用いた可視化グラフが望ましいのかを客観的に評価可能となり, 
検証実験では, 定性評価での結果と一致することを確認し, 提案法の妥当性を検証した. 
今後の研究では，提案法の妥当性を様々な可視化手法で検証する．

\appendix
\chapter{可視化ブラウジングツールの利用方法}

実際に, アップロードしたクエリ文書から, 検索結果文書を表示する可視化ブラウジングツールを開発した.
当ツールの機能説明とシステム構成図を明記することで, ツール全体の設計概要を示す.

\section{クエリ文書アップロード}

当ツールは, 図~\ref{fig7-1}(a)に示す初期画面から, 
画面上部に位置する黄緑色のヘッダーを押して任意のクエリ文書をアップロードし,
「Submit」ボタンを押すことで, 検索結果類似文書を図~\ref{fig7-1}(b)のようにグラフ可視化できる仕組みである.
なお, グラフ可視化するまでに多少の待機時間が発生するため, 
その間は図~\ref{fig7-2}に示すロード画面を表示する.

\begin{figure}[H]
   \begin{center}
    \includegraphics[width=\hsize]{figs/ku/fig7-1/fig1a.jpg} 
    \mbox{(a) 初期画面}
   \end{center}
   \begin{center}
    \includegraphics[width=\hsize]{figs/ku/fig7-1/fig1b.jpg}
    \mbox{(b) 可視化グラフ生成後の画面}
   \end{center}
 \caption{可視化ブラウジングツール}
 \label{fig7-1}
\end{figure}
\begin{figure}[H]
  \begin{center}
  \includegraphics[width=\hsize]{figs/ku/fig7-2/fig1.jpg}
  \end{center}
  \caption{可視化グラフ生成中のロード画面}
  \label{fig7-2}
\end{figure}

\section{ラジオボタン設定}

図~\ref{fig7-1}(a),(b) に示す, 画面上部に配置したラジオボタンの集合で構成される３つの設定項目 
「Visual, Graph, Search」は, それぞれ,
\begin{enumerate}
  \item グラフ可視化手法
  \item グラフ構成法
  \item 検索文書数
\end{enumerate}
を指定するために用意している.

\section{カテゴリー可視化}

可視化グラフの上部に配置している「● 説明語：X」は,
各カテゴリーに属するノード数を示している.
図~\ref{fig7-1}(b)の場合は, 「iphone」という説明語を持つノード数が39個で, 
「カバー」が30個, 「アプリ」は8個存在することを示しており,
その他の組も同様の意味を表している.
なお, 全てのカテゴリーを列挙している訳ではなく, 属するノード数が３以上であるカテゴリーのみを表示している.
これは, 無数のカテゴリーが列挙されると, 返って視認性が悪くなり, 
結果的に検索結果文書の全体像把握に時間が掛かってしまう問題を考慮した故の策である.

\section{ノード}

当ツールでは, 任意のノード(文書)をホバーすることで, 
図~\ref{fig7-3}(a)に示すような, 下記項目を閲覧できる.
\begin{enumerate}
  \item 説明語
  \item 文書ID
  \item 文書名
  \item 文書内容
\end{enumerate}

さらに, クリックすることで, 該当ノード(文書)の文書内容を,
図~\ref{fig7-3}(b)に示す文書詳細閲覧画面の右側に全て表示する.
このように配置することで, クエリ文書と任意の検索文書を容易に比較できる.

\begin{figure}[H]
   \begin{center}
    \includegraphics[width=\hsize]{figs/ku/fig7-3/fig1a.jpg}
    \mbox{(a) ホバー時に表示される文書概要タイル}
   \end{center}
   \begin{center}
    \includegraphics[width=\hsize]{figs/ku/fig7-3/fig1b.jpg}
    \mbox{(b) クリック時に表示される文書内容詳細画面}
   \end{center}
 \caption{ノードを対象とした動作確認}
 \label{fig7-3}
\end{figure}

\section{システム構成}

当ツールのシステム構成図を図~\ref{fig7-4}に示す.
可視化グラフを描画するクライアント側では JavaScript を用い,
各画面コンポーネントのレイアウト構成には formantic-UI を使用した.
なお, アップロードされるクエリ文書の形態素解析がなされていない可能性を考慮し,
事前にMeCab~\cite{kudo04}で前処理を実行してからサーバーサイド側へリクエストを送信する.
これら環境は Docker 上で構築した.
また, 可視化グラフの描画に必要なデータを生成するサーバーサイド側の実装に関しては, 
既存の研究で用いられていたC言語のコードを, 本研究の趣旨に適する処理へ改良し, 
Apache上に構築した CGI（Common Gateway Interface）を仲介する形で, 
グラフ可視化に必要なデータをクライアント側へレスポンス可能な構成にした.

\begin{figure}[H]
  \begin{center}
  \includegraphics[width=\hsize]{figs/ku/fig7-4/fig1.jpg}
  \end{center}
  \caption{可視化ブラウジングツールのシステム構成図}
  \label{fig7-4}
\end{figure}

%=====================================================================================
\chapter*{謝辞} 
\addcontentsline{toc}{chapter}{謝辞} 
本論文を作成するにあたり，指導教官として終始多大なご指導を賜った，斉藤和巳教授に感謝の意を表します．

%=====================================================================================
\renewcommand{\bibname}{参考文献}
\begin{thebibliography}{10}
\bibitem{ozaki1}
尾崎 了祐, 斉藤和己.
類似貢献度に基づく文書検索結果の可視化分析法.
情報処理学会第８６回全国大会, 2024. 
\bibitem{ozaki2}
尾崎 了祐, 斉藤和己.
類似貢献度に基づく文書検索結果の可視化評価法.
情報処理学会第８７回全国大会, 2025. 
\bibitem{saito}
斉藤和巳.
ネットワークの可視化技術―大規模情報からの意味情報の抽出―.
数理科学, 536, 78-83, 2008.
\bibitem{ikeda17}
池田 弘行.
機械学習を用いた 文章データの解析・可視化技術.
東芝レビュー, 5, 68-69, 2017. 
\bibitem{oda04}
小田 良治, 土井 章男.
大規模文書検索結果のクラスタリングと可視化.
情報処理学会第66回全国大会, 2004. 
\bibitem{kudo04}
T. Kudo, K. Yamamoto, and Y. Matsumoto, 
Applying Conditional Random Fields to Japanese Morphological Analysis, 
Proceedings of the 2004 Conference on Empirical Methods in Natural Language Processing (EMNLP-2004), 230-237, 2004.
\bibitem{kamada89}
T. Kamada and S. Kawai. 
An algorithm fordrawing general undirected graphs. 
Information Processing Letters. 31, 7-15, 1989.
\bibitem{BETT}
G. Battista, P. Eades, R. Tamassia, and I. Tollis, 
Graph drawing: An annotated bibliography Prentice-Hall, New Jersy, 1999.
\bibitem{C}
F. R. K. Chung,
Spectral Graph Theory, Number 92 in CBMS Regional Conference Series 
in Mathematics American Mathematical Society, 1997. 
\bibitem{YST}
山田 武士, 斉藤 和巳, 上田 修功, クロスエントロピー最小化に基づくネットワークデータの埋め込み.
情報処理学会論文誌, 44 2401-2408, 2003.

\end{thebibliography}

\end{document}
